\section{The natural octahedral pair: a sharp uniform angle}
\label{sec:octahedral-pair}

We construct a pair of octahedral subgroups whose fixed-space angle is
controlled by the Ramanujan bound for a quaternionic Cayley walk.
A compression identity connects the two operators. Throughout, $q$ is a prime,
$q\equiv\pm3\pmod8$, $q\ne3$, and $G_q=\PGL_2(\F_q)$.
When a representation $\rho$ is fixed, we write $P_H=P_H^\rho$;
we retain the superscript when taking maxima over representations.

\subsection{The pair and its four neighbours}
Let $\mathbb H=(-1,-1)_{\mathbb Q}$, with $i^2=j^2=k^2=ijk=-1$,
and take projective classes in $\mathbb H^\times/\mathbb Q^\times$:
\[
 A_0=\langle i,1+i,1+i+j+k\rangle,\qquad h_0=i+j+k,
 \qquad S_0=\{\pm i\pm j\pm k\}/\{\pm1\}.
\]
The four elements of $S_0$ have norm $3$ and are projective involutions.
For any odd prime $q$, the quaternion algebra over $\F_q$ is split.
Indeed, the two sets of squares and of $-1$ minus squares have more
than $q$ elements in total, so $x^2+y^2=-1$ has a solution; this gives
an explicit splitting into $M_2(\F_q)$~\cite[Proposition 2.5.2]{DSV2003}.
Reducing the indicated projective elements defines $A,h,S$ in $G_q$.
Their norms $1,2,4,3$ are invertible for our primes. Changing the splitting
conjugates the entire construction, leaving the norms below unchanged.

\begin{lemma}[Octahedral reduction]
\label{lem:octahedral-construction}
One has $A\cong S_4$, with exterior transpositions, and
\[
 B=hAh^{-1},\quad A\cap B\cong S_3,\quad h^2=1,\quad
 \langle A,B\rangle=G_q,\quad AhA/A=\{sA:s\in S\}.
\]
The last set has four distinct cosets, and $S$ is stable under conjugation
by $A$.
\end{lemma}
\begin{proof}
On the trace-zero subspace with basis $(i,j,k)$, conjugation by the three
generators of $A_0$ and by $h_0$ gives, respectively,
\[
 \begin{pmatrix}1&0&0\\0&-1&0\\0&0&-1\end{pmatrix},\quad
 \begin{pmatrix}1&0&0\\0&0&-1\\0&1&0\end{pmatrix},\quad
 \begin{pmatrix}0&0&1\\1&0&0\\0&1&0\end{pmatrix},\quad
 H=\frac13\begin{pmatrix}-1&2&2\\2&-1&2\\2&2&-1\end{pmatrix}.
\]
The first three generate the $24$ determinant-one signed permutation
matrices. Their action on the four body diagonals identifies them with
$S_4$; they remain distinct modulo every odd prime. The adjoint action
is faithful projectively, also in odd characteristic. Conjugation by
$\alpha\in\mathbb H(\F_q)^\times$ is $L_\alpha R_{\alpha^{-1}}$ on
$\mathbb H(\F_q)$, of determinant $N(\alpha)^2N(\alpha)^{-2}=1$ for the
norm $N$; it fixes $1$ and preserves the norm form, so the adjoint action
embeds $G_q$ into $\mathrm{SO}_3(\F_q)$. The norm of $1+i$
is $2$, a nonsquare modulo $q$, so the determinant-square character is
non-trivial on $A$ and hence is its sign character. In particular the
transpositions of $A$ are exterior.

Here is a uniform intersection calculation. Let $K$ be the stabilizer in
$A$ of the line through $(1,1,1)$: its matrices are the positive even
permutation matrices and the negative odd permutation matrices. Thus
$K\cong S_3$, and every element of $K$ commutes with $H$.
Let $V$ be the four diagonal sign matrices of determinant one. Then
$A=VK$, and $V\cap K=1$. The three non-identity elements of $V$ are
conjugate under $K$. For $d=\operatorname{diag}(1,-1,-1)$,
\[
 HdH=\frac19\begin{pmatrix}-7&-4&-4\\-4&-1&8\\-4&8&-1\end{pmatrix}.
\]
Its second row has three nonzero entries modulo every prime $q\ge5$,
so $HdH\notin A$. Since $H$ commutes with $K$, the same conclusion
holds for the other two non-identity cosets of $K$ in $A$.
Therefore $A\cap HAH=K$, in characteristic zero and modulo every
allowed prime. This also proves $h\notin A$ and $B\ne A$.
The exterior $S_4$ is maximal in $G_q$ for these primes by
\cite[Theorem 3.5]{Giudici2007}; hence $\langle A,B\rangle=G_q$.

Finally $H^2=I$, and $A$ acts transitively by conjugation on the four
body diagonals $S$, with stabilizer $K$. The map
$A/K\longrightarrow AhA/A$, $aK\mapsto ahA$, is bijective by the
intersection calculation. Since $(aha^{-1})A=ahA$, its images are
exactly the four cosets represented by $S$.
\end{proof}

\begin{proposition}[The intersection-six stratum]
\label{prop:octahedral-cap6}
Let $A_1\ne A_2$ be exterior octahedral subgroups of $G_q$. Then
$|A_1\cap A_2|\le6$. If $|A_1\cap A_2|=6$, the ordered pair $(A_1,A_2)$
is simultaneously conjugate in $G_q$ to $(A,hAh^{-1})$, and for fixed
$A_1$ there are exactly four such $A_2$.
\end{proposition}
\begin{proof}
The exterior octahedral subgroups form one conjugacy class (Dickson's
classification of the subgroups of $\PSL_2(q)$ and
$\PGL_2(q)$~\cite{Huppert1967,King2005}; for the statement in this form
see~\cite[proof of Proposition~3.4]{Giudici2007}), so we may take $A_1=A$ and
$A_2=gAg^{-1}$ with $g\notin A$, in the orthogonal model of
\Cref{lem:octahedral-construction}.

\emph{Normalisers.} The three non-identity elements of $V$ have the
coordinate axes as fixed lines, so an element normalising $V$ permutes
these axes. Orthogonality forces its nonzero entries to be $\pm1$, and
determinant one leaves the $24$ elements of $A$. Thus $N_{G_q}(V)=A$;
since $V$ is characteristic in $A$, also $N_{G_q}(A)=A$.

\emph{Orders $8$ and $12$.} Put $D=A\cap A_2$. The determinant-square
character restricts to the sign character of $A$ and of $A_2$, so
$A\cap\PSL_2(q)$ and $A_2\cap\PSL_2(q)$ are their $A_4$ subgroups, whose
only subgroups of order $4$ are their normal Klein four subgroups. If
$|D|=12$, then $D$ is the $A_4$ of both groups. If $|D|=8$, then $D$
contains transpositions, so $D\cap\PSL_2(q)$ has order $4$ and is the
normal Klein four of both groups. In either case $A_2$ normalises $V$,
so $A_2\le A$, a contradiction. Hence $|D|\le6$, and $|D|=6$ forces
$D\cong S_3$.

\emph{Order $6$.} The four $S_3$ subgroups of $A$ are conjugate in $A$,
so replacing $g$ by $a_1ga_2^{-1}$ with $a_1,a_2\in A$, which does not
change $AgA$, we may assume that $g$ normalises $K$. Every automorphism of
$S_3$ is inner, so $N_{G_q}(K)=K\,C_{G_q}(K)$. The elements of $K$ are the
matrices $\operatorname{sgn}(\pi)P_\pi$, $\pi\in S_3$, so a matrix commutes
with $K$ exactly when it commutes with every permutation matrix $P_\pi$,
that is, when it is constant on the two orbits of $S_3$ on index pairs:
the commutant is $\{aI+bJ\}$, with $J$ the all-ones matrix. For
$M=aI+bJ$, orthogonality gives $a^2=1$ and $b(2a+3b)=0$; since $q\ne3$,
$M$ acts by $a+3b$ on $(1,1,1)$ and by $a$ on the complementary plane
$x+y+z=0$, so determinant one gives $a+3b=1$. The solutions are $I$ and
$H$. Hence $C_{G_q}(K)=\{1,h\}$, $N_{G_q}(K)=K\sqcup Kh$, and $g\in Kh$
because $g\notin A$. Therefore $AgA=AhA$; writing $g=ahb$ with $a,b\in A$
gives $A_2=a(hAh^{-1})a^{-1}$, so $(A,A_2)=a(A,hAh^{-1})a^{-1}$.

\emph{Count.} As $N_{G_q}(A)=A$, the map $gA\mapsto gAg^{-1}$ is
injective. The subgroups $A_2$ with $|A\cap A_2|=6$ are the images of the
cosets in $AhA/A$, of which there are $[A:K]=4$ by
\Cref{lem:octahedral-construction}.
\end{proof}

\subsection{Compression and the arithmetic spectral input}
Put $\mu=\frac14\sum_{s\in S}\delta_s$ and let $e_A$ denote the
uniform average over $A$ in $\C[G_q]$.

\begin{lemma}[The compression identity]
\label{lem:octahedral-compression}
In the group algebra,
\begin{equation}\label{eq:octahedral-compression}
 \mu e_A=e_A\mu=e_Ahe_A.
\end{equation}
Consequently $\rho(\mu)$ commutes with $P_A$ and its restriction to
$V_\rho^A$ is $P_A\rho(h)P_A$, whose norm equals $\opn{P_AP_B}$.
\end{lemma}
\begin{proof}
Conjugation invariance of $S$ gives commutation with $A$. Averaging
$ahe_A$ over $a\in A$ visits each of the four cosets in
\Cref{lem:octahedral-construction} equally often, proving
$e_Ahe_A=\frac14\sum_{s\in S}se_A$. Apply $\rho$ and
\Cref{cor:double-coset-angle}.
\end{proof}

\begin{proposition}[Quaternionic Ramanujan input]
\label{prop:octahedral-arithmetic}
For $\chi_q(g)$ the determinant-square character of $G_q$ and
$\epsilon_q=(3/q)$, the operator $\rho(\mu)$ has norm at most
$\sqrt3/2$ on every irreducible $\rho\notin\{\mathbf1,\chi_q\}$.
On those two characters its values are $1$ and $\epsilon_q$.
\end{proposition}
\begin{proof}
We apply the theorem of Davidoff--Sarnak--Valette
\cite[Theorem 4.2.2]{DSV2003}. Its hypotheses are distinct odd primes
$p,q$ with $q>2\sqrt p$, and its distinguished norm-$p$ generators
are specified in Section~2.6 of that reference. For $p=3$ they are
exactly the four classes of $\pm i\pm j\pm k$, so every prime here
satisfies the hypotheses. Remark~4.4.7 of the same reference derives the
Ramanujan bound from Eichler's theorem on weight-two cusp forms.

The vertex set of the walk follows from the construction. The set $S$ is
stable under conjugation by $A$ and contains $h$, so $\langle S\rangle$ is
normalised by $\langle A,h\rangle=G_q$ and is a non-trivial normal subgroup
of $\PGL_2(q)$. Since $\PSL_2(q)$ is simple with trivial centraliser in
$\PGL_2(q)$, $\langle S\rangle$ contains $\PSL_2(q)$. Every $s\in S$ has
determinant class $3$; hence $\langle S\rangle=\PSL_2(q)$ when
$\epsilon_q=1$ and $\langle S\rangle=\PGL_2(q)$ when $\epsilon_q=-1$. These
are the two vertex sets in the cited theorem, and the external input is
the spectral bound alone.

If $\epsilon_q=-1$, their Cayley graph on $G_q$ is connected,
bipartite and $4$-regular, with all eigenvalues other than $\pm4$
bounded in absolute value by $2\sqrt3$. The exceptional normalized
eigenspaces are the constants and $\chi_q$.
If $\epsilon_q=1$, the theorem gives a connected nonbipartite
Ramanujan graph on $\PSL_2(q)$. On $G_q$ the same walk is the union
of two copies, so the exceptional eigenspace at $1$ is spanned by
the constants and $\chi_q$, and there is no eigenvalue $-1$.
The regular representation contains every irreducible, giving the
assertion. In both cases each $s\in S$ has determinant square class $3$.
\end{proof}

The compression identity \eqref{eq:octahedral-compression} identifies the
Hecke operator. Left multiplication
by the fixed set $S$ on the cosets $G_q/A$ is generally a different
operator from the Hecke adjacency. For $q=5$, those two unnormalized
spectra are respectively $\{4,3,3,3,-1\}$ and
$\{4,-1,-1,-1,-1\}$. The latter graph is $K_5$.

\begin{theorem}[Sharp angle for the natural pair]
\label{thm:octahedral-sharp}
For the pair in \Cref{lem:octahedral-construction}, set
\[
 c_q=\max_{\rho\in\widehat{G_q},\,\rho\ne\mathbf1}
          \opn{P_A^\rho\rho(h)P_A^\rho}.
\]
Then, for every pair $A_1,A_2$ of exterior octahedral subgroups of $G_q$
with $|A_1\cap A_2|=6$,
\[
 \max_{\rho\ne\mathbf1}\opn{P_{A_1}^\rho P_{A_2}^\rho}
 =c_q=\max_{\rho\ne\mathbf1}\opn{P_A^\rho P_B^\rho}\le\frac{\sqrt3}{2},
 \qquad
 \lim_{q\to\infty}c_q
       =\frac{\sqrt3}{2}.
\]
Thus the constant is best possible uniformly over this family.
\end{theorem}
\begin{proof}
By \Cref{prop:octahedral-cap6}, $(A_1,A_2)=x(A,B)x^{-1}$ for some
$x\in G_q$, and $P_{xAx^{-1}}^\rho=\rho(x)P_A^\rho\rho(x)^{-1}$; this
gives the first equality in every representation. It remains to treat
$(A,B)$. For $\rho\notin\{\mathbf1,\chi_q\}$, the compression identity and
\Cref{prop:octahedral-arithmetic} give the upper bound. On $\chi_q$,
$P_A=0$ because $A$ is exterior, so this character contributes zero.

For the lower bound, define $X_q$ with vertices the left cosets $xA$,
and neighbours $xsA$ for $s\in S$. This is independent of the
representative $x$, since $S$ is invariant under conjugation by $A$.
The four neighbours are distinct and none is $xA$, by
\Cref{lem:octahedral-construction}. The relation is symmetric because
$h^{-1}=h$ and $AhA=(AhA)^{-1}$. The graph is connected because
$\langle A,h\rangle=G_q$, and has $\abs{G_q}/24$ vertices.
On functions on $G_q$ invariant under right multiplication by $A$,
its normalized adjacency is the restriction of the \emph{right}
convolution
\[
 (Wf)(x)=\frac14\sum_{s\in S}f(xs).
\]
Under the finite Peter--Weyl decomposition, the multiplicity spaces of
this coset representation are $V_\rho^A$ (or their duals, with identical
norms), and the blocks are the compressions
$P_A\rho(h)P_A$. Hence the absolute nonconstant spectral radius of
$X_q$ is $4c_q$. The trivial representation contributes just the
constant functions, and $\chi_q$ has no $A$-fixed vectors.
Alon--Boppana~\cite[Theorem 1.3.1]{DSV2003}, applied to these connected
loopless $4$-regular graphs with vertex count tending to infinity,
gives $\liminf 4c_q\ge2\sqrt3$. Together with the upper bound this
proves the limit.
\end{proof}

\subsection{From angles to tree blocks}
The minimum fixed-space angle $\theta_q=\arccos c_q$ satisfies
$\theta_q\ge\pi/6$ and $\theta_q\to\pi/6$ by
\Cref{thm:octahedral-sharp}. Since
$V_\rho^A\cap V_\rho^B=V_\rho^{G_q}=0$ for every non-trivial irreducible,
these are Friedrichs angles, with angle $\pi/2$ when a fixed space is zero.
This angular separation controls how two subgroup averages interact.
The following estimate transfers it to tree blocks that
approximate those averages.

\begin{lemma}[Quantitative gluing]
\label{lem:octahedral-gluing}
Let $P,Q$ be orthogonal projections with $\opn{PQ}\le c<1$. Suppose
$X=P+R$, $Y=Q+S$, where
$PR=RP=QS=SQ=0$ and $\opn R,\opn S\le\lambda<1$. Then
\begin{equation}\label{eq:octahedral-gluing}
 \opn{XYX^\dagger Y^\dagger}^2
 \le1-\frac{(1-\lambda)(1-c)}{8(1+\lambda)}.
\end{equation}
\end{lemma}
\begin{proof}
Put $g=1-\lambda^2$. Each factor $Z$ with associated projection $F$
satisfies
\[
 \norm w^2-\norm{Zw}^2\ge g\norm{(I-F)w}^2,
 \qquad \norm{w-Zw}\le(1+\lambda)\norm{(I-F)w}.
\]
Apply the four factors in their actual order to a unit $v$, and denote
their norm-squared losses by $e_j\ge0$. Their sum is
$E=1-\norm{XYX^\dagger Y^\dagger v}^2$.
At each step, telescoping back to $v$ and Cauchy--Schwarz give
\[
 \norm{(I-F)v}\le\frac{1+\lambda}{\sqrt g}
             \sum_{i\le j}\sqrt{e_i}
       \le\frac{2(1+\lambda)}{\sqrt g}\sqrt E.
\]
Both $P$ and $Q$ occur. Also $\opn{P+Q}\le1+c$: the map
$v\mapsto(Pv,Qv)$ has squared norm bounded by that of the block matrix
with identity diagonal and off-diagonal blocks $PQ,QP$.
Consequently
\[
 1-c\le\norm{(I-P)v}^2+\norm{(I-Q)v}^2
       \le\frac{8(1+\lambda)^2}{g}E.
\]
Taking the supremum over $v$ proves \eqref{eq:octahedral-gluing}.
\end{proof}

\begin{corollary}[Uniform contraction for natural-pair blocks]
\label{cor:octahedral-block}
Suppose a valid level-$3$ tree $\tau$ has level-$2$ leaf subgroups
$A,B\cong S_4$, exterior with $|A\cap B|=6$, and all leaves are exterior
involutions. For every non-trivial $\rho$,
\[
 \opn{M_\tau(\rho)}\le\kappa,
 \qquad
 \kappa=\left(1-\frac{183}{1736}
                       \left(1-\frac{\sqrt3}{2}\right)\right)^{1/2}<1.
\]
For a valid level-$m$ tree, $m\ge3$, all of whose level-$3$ sub-blocks
have this property, replace $\lambda$ by $\kappa$ and $D$ by $3$ in
\Cref{prop:quant}.
\end{corollary}
\begin{proof}
On the fixed space of its leaf subgroup, each level-$2$ operator is the
identity. On the orthogonal complement its norm is at most
$\lambda=17/200$: this is a finite rational certificate on the four
non-trivial irreducibles of $S_4$, over its $432$ valid ordered
transposition configurations generating $S_4$. The exact certificate
described in \Cref{sec:computation} constructs rational matrices with
invariant Gram matrices $G_\pi$ and verifies all principal minors of
$(17/200)^2G_\pi-D^tG_\pi D$ are nonnegative. Restriction of any ambient
unitary representation to $S_4$ decomposes into these irreducibles and
the trivial one; hence the certificate applies in every ambient $G_q$.
Both invariant subspaces are preserved, so the hypotheses of
\Cref{lem:octahedral-gluing} hold. Insert $c=\sqrt3/2$ from
\Cref{thm:octahedral-sharp}. The last assertion follows by expanding
into the level-$3$ blocks, as in \Cref{prop:quant}.
\end{proof}

The sharp angle estimate and the finite $S_4$ certificate thus give the
explicit contraction bound $\kappa$ for these octahedral blocks. Applying
\Cref{prop:stability} inside $S_4$ also gives a purely analytic, weaker
constant in the same corollary.
